\scanpage{184}\mspage{39}
on arrive à réconstituer le type d'homotopie via son $\infty$-groupoïde fondamental, comme expliqué dans ma lettre précédente, \deleted{\unsure{e|a}}et rappelé dans celle-ci. (Cela signifie en particulier que toutes les opérations cohomologiques possibles et imaginables sont déjà incluses dans la donnée d'un tel système de $n$-catégories .\,.\,.\,.) \deleted{Mais}

De même, les types d'homotopie plus \unsure{cannelés|canelés}, liés à certains diagrammes finis, qu'on peut associer à \deleted{\unsure{des}}certains types de stratifications \inserted{$\mathcal{P}$} d'espaces topologiques ``modérés'' $X$ disons, devraient correspondre de façon aussi parfaite aux \deleted{\unsure{$n$-}}\inserted{$(n+1)$-}catégories de $n$-champs sur $X$ qui sont loc constant sur chacune des strates de $\mathcal{P}$ (disons: qui sont subordonnés à $\mathcal{P}$). Si plus haut le \unsure{topos|topis} de la cat. dérivée loc. constante marchait bel et bien, on s'attendrait ici à récupérer le type d'homotopie mixte à partir de la \inserted{ss-}catégorie de la catégori dérivée des faisceaux abéliens sur $X$, formée des complexes qui sont à \unsure{coh.|ch.} loc. constant sur chacune des strates -- avec en plus les opérations $\overset{\mathbb{L}}{\otimes}$, $\mathbb{R}\underline{\mathrm{Hom}}$ plus, les
% archive sheet number: 184
